\chapter{Introduzione}
\label{cap:introduzione}

\section{Problema di ricerca}
\label{sec:problema}

Nel \emph{quantum machine learning} ibrido un circuito parametrizzato opera come
componente di una pipeline classica, e i suoi parametri vengono ottimizzati
insieme a quelli della rete \parencite{biamonte2017qml,cerezo2021variational}. In
questa tesi il circuito, indicato come \emph{Variational Quantum Circuit} (VQC),
riceve un vettore di attivazioni classiche, lo codifica nello stato di un registro
di qubit, lo trasforma con una sequenza di porte i cui angoli sono parametri
addestrabili e restituisce alla rete i valori misurati. L'interfaccia di
autodifferenziazione lo presenta all'ottimizzatore come un modulo qualsiasi, con
pesi, output e gradiente \parencite{bergholm2018pennylane}.

Questa interfaccia consente di inserire un VQC al posto di uno strato lineare,
purché le dimensioni siano compatibili, ma non stabilisce se la sostituzione sia
utile. La domanda è quindi empirica:
si sostituisce lo strato, si tiene fermo tutto il resto, si misura.

\subsection{Quattro promesse distinte}
\label{ssec:quattro-promesse}

In questo lavoro l'espressione «vantaggio quantistico» viene scomposta in quattro
proposizioni verificabili separatamente. La distinzione evita di trasferire
all'accuratezza conclusioni che riguardano, per esempio, la complessità
computazionale o la geometria del kernel:

\begin{enumerate}[label=(\alph*)]
  \item \textbf{Speedup computazionale.} Il calcolo eseguito dal circuito è
    intrattabile classicamente, o comunque asintoticamente più costoso da
    simulare che da eseguire su hardware quantistico.
  \item \textbf{Espressività.} A parità di numero di parametri addestrabili, la
    famiglia di funzioni realizzabile da un VQC è più ricca di quella di uno
    strato classico di pari dimensione, e questo si traduce in un errore di
    approssimazione inferiore.
  \item \textbf{Regolarizzazione implicita.} L'unitarietà dell'evoluzione e la
    limitatezza dei valori di aspettazione vincolano la rappresentazione prodotta
    dal circuito; il vincolo agisce come un regolarizzatore che non compare nella
    funzione di costo e riduce l'\emph{overfitting}.
  \item \textbf{Separabilità.} La mappa di codifica proietta i dati in uno spazio
    di Hilbert di dimensione $2^{\nq}$; il prodotto interno fra stati codificati
    definisce un kernel che potrebbe separare le classi meglio di un kernel
    classico \parencite{schuld2019quantum,havlicek2019supervised}.
\end{enumerate}

Questa tesi affronta soltanto le affermazioni (b), (c) e (d). L'affermazione (a)
non è valutabile con il protocollo adottato: tutti gli esperimenti
qui descritti girano sul simulatore di vettore di stato \code{default.qubit} della
libreria PennyLane \parencite{bergholm2018pennylane}, che rappresenta in memoria
le $2^{\nq}$ ampiezze e le fa evolvere con prodotti matrice-vettore. Con i valori
di $\nq$ impiegati in questo lavoro, fra 4 e 10 qubit, la simulazione è
deterministica, priva di rumore di campionamento e sostenibile sulle risorse
disponibili. Nessuna misura di tempo riportata nel
\cref{cap:risultati} può quindi essere letta come evidenza a favore o contro uno
speedup.

\subsection{La domanda specifica}
\label{ssec:domanda}

Ristretta alle tre affermazioni verificabili in simulazione, la domanda di ricerca
è la seguente:

\begin{quote}
\itshape
Negli esperimenti considerati, l'inserimento di un circuito quantistico
variazionale in un Transformer testuale o visivo, oppure l'uso di un kernel di
fedeltà su immagini mediche, produce un miglioramento misurabile della qualità
predittiva, del gap di generalizzazione o della separabilità rispetto a controlli
classici appaiati?
\end{quote}

Il Transformer \parencite{vaswani2017attention} offre punti di intervento
circoscritti, e il Vision Transformer trasferisce la stessa struttura alle
sequenze di patch di immagine \parencite{dosovitskiy2020image}. In entrambe le
architetture esistono punti di sostituzione ben identificati: le tre proiezioni
$Q$, $K$, $V$ di ogni testa di attenzione e la proiezione delle patch. Sono state
proposte architetture quantistiche che intervengono sull'attenzione o su altri
blocchi del ViT \parencite{cherrat2024quantumvision,comajoan2024qvt}, e questa
tesi non tenta di riprodurle integralmente: adotta sostituzioni più piccole per
appaiare inizializzazione, dati e metriche di addestramento e misurare la
variabilità fra più seed. Il contributo metodologico è il confronto appaiato tra
le varianti quantistiche e le rispettive baseline classiche.

\section{Ipotesi}
\label{sec:ipotesi}

Il lavoro è organizzato attorno a quattro ipotesi operative. Ciascuna è enunciata
con una metrica primaria, una direzione attesa e un criterio esplicito che ne
determina l'esito.

\begin{table}[htbp]
\centering
\caption[Le quattro ipotesi e i loro criteri di decisione]{%
Le quattro ipotesi operative. Ogni riga trova risposta in una riga della tabella
riassuntiva della \cref{sec:verifica-ipotesi}. I pedici $q$ e $c$ indicano
rispettivamente la variante quantistica e la baseline classica appaiata.}
\label{tab:ipotesi}
\small
\begin{tabularx}{\textwidth}{@{}c X l l@{}}
\toprule
\textbf{\#} & \textbf{Enunciato} & \textbf{Metrica primaria} & \textbf{Direzione attesa} \\
\midrule
H1 & Le proiezioni $Q,K,V$ quantistiche migliorano la modellazione del linguaggio
   & validation loss & $\mathcal{L}_q < \mathcal{L}_c$ \\[4pt]
H2 & Il \emph{patch embedding} quantistico migliora la classificazione di immagini
   & test accuracy & $\accte^{\,q} > \accte^{\,c}$ \\[4pt]
H3 & Il VQC agisce da regolarizzatore implicito
   & gap $\gap = \acctr - \accte$ & $\gap_{\text{quantum}} < \gap_{\text{bounded}}$ \\[4pt]
H4 & Il kernel di fedeltà quantistica separa le classi meglio della similarità coseno
   & criterio di Fisher & $\text{Fisher}_q > \text{Fisher}_c$ \\
\bottomrule
\end{tabularx}
\end{table}

Il criterio di decisione è comune a tutte e quattro ed è formulato sulle differenze
appaiate $d_s = m_{q,s} - m_{c,s}$ calcolate sullo stesso seed $s$:

\begin{quote}
Un'ipotesi operativa è \textbf{supportata} se l'intervallo di confidenza al
\SI{95}{\percent} della differenza media $\dmean$ ricade interamente nella
direzione attesa, ed è \nonsupportata{} se contiene lo zero. Se l'intervallo
ricade interamente nella direzione opposta, l'esito è registrato come
\evidenzacontraria. Questa convenzione riguarda il
sostegno empirico all'ipotesi di ricerca e non va confusa con il rifiuto o il
mancato rifiuto dell'ipotesi nulla nel test statistico.
\end{quote}

Per H3 il criterio è applicato ai tre dataset, tenendo conto dei confronti
multipli come descritto nel \cref{cap:metodologia}. L'ablazione del
\cref{cap:esperimenti} costruisce tre varianti di \emph{patch embedding}: una
lineare non vincolata (\code{vanilla}), una classica con output limitato in
$[-1,1]$ (\code{bounded\_mlp}) e una quantistica (\code{quantum\_reg}). Il
confronto primario è con \code{bounded\_mlp}, perché entrambe le varianti hanno
output limitato. MLP e VQC appartengono comunque a famiglie funzionali diverse;
il confronto non isola una singola proprietà quantistica
(\cref{sec:esp-ablazione}).

Nessuna delle quattro ipotesi operative risulta supportata nella direzione
attesa. Questo esito riguarda le configurazioni esaminate e non costituisce un
giudizio sulla letteratura richiamata nella \cref{ssec:quattro-promesse}.

\section{Contributi}
\label{sec:contributi}

\paragraph{Contributi sperimentali.} Quattro esperimenti distinti, uno per
ipotesi, ciascuno con baseline classica appaiata e protocollo multi-seed; per
l'esperimento sulla regolarizzazione, anche multi-dataset. L'analisi impiega
differenze appaiate, intervalli di confidenza parametrici e bootstrap, test di
Wilcoxon dei ranghi con segno e dimensione dell'effetto $\dz$, riportati sempre
insieme e mai ridotti al solo valore $p$.

\paragraph{Contributi software.} La libreria \code{quantum\_framework}, descritta
nel \cref{cap:software}, è stata estratta per rifattorizzazione dai quattro
esperimenti, originariamente sviluppati come repository separati. Contiene i due
layer variazionali riutilizzabili, il caricamento dei dati MedMNIST con
generazione deterministica di corruzioni, un ciclo di addestramento con
\emph{hook} diagnostici su gradienti e saturazione delle attivazioni, le metriche
di classificazione, le funzioni di confronto statistico e la raccolta dei metadati
di provenienza. A corredo: un entry point unico (\code{run.py}), uno schema comune
dei risultati versionato e, per la run più costosa, la possibilità di
interrompere e riprendere l'esecuzione.

\paragraph{Contributo metodologico.} Il design adotta baseline appaiate.
Nell'ablazione, la baseline classica ha più parametri della controparte quantistica nel
\emph{patch embedding} (\num{1880} contro \num{1540}); negli esperimenti 01 e 02
è invece la variante quantistica ad averne rispettivamente 336 e 48 in più. I
conteggi escludono uno svantaggio sistematico della variante quantistica nel
numero di parametri, ma non dimostrano parità di capacità funzionale.

L'infrastruttura sviluppata rende ripetibile il confronto appaiato; la libreria
di simulazione quantistica è invece una dipendenza del progetto.

\section{Struttura della tesi}
\label{sec:struttura}

I \cref{cap:fondamenti-quantistici,cap:transformer,cap:vit} introducono i concetti
quantistici, il Transformer e il Vision Transformer, fino ai punti di inserimento
dei circuiti. Il \cref{cap:software} descrive l'infrastruttura e il
\cref{cap:metodologia} il protocollo sperimentale con le minacce alla validità.

Il \cref{cap:esperimenti} descrive il design dei quattro esperimenti, il
\cref{cap:risultati} ne riporta i risultati e il \cref{cap:discussione} li
interpreta. Il
\cref{cap:conclusioni} risponde alle quattro ipotesi e indica gli sviluppi
successivi. Le appendici raccolgono configurazioni, comandi di riproduzione,
tabelle per singolo seed, diagrammi dei circuiti e l'inventario delle esecuzioni.
